\section{Equilibrium Analysis}

\begin{frame}
  \frametitle{Three Operating Modes}
  \footnotesize
  This section first defines the three operating modes available to $M$,
  then derives the equilibrium outcome under each mode.

  \vspace{0.15cm}
  \begin{table}
    \centering
    \scriptsize
    \renewcommand{\arraystretch}{1.35}
    \begin{tabularx}{\textwidth}{@{}p{2.35cm}X X X@{}}
      \toprule
      \textbf{Mode chosen by $M$}
        & \textbf{Platform $M$}
        & \textbf{Seller $S$}
        & \textbf{Fringe sellers} \\
      \midrule
      Marketplace mode
        & Hosts the marketplace and sets commission $\tau$; does not sell its own product.
        & Joins the marketplace, chooses innovation $\Delta$, inside price $p_i$, and outside price $p_o$.
        & Provide the competitive benchmark on and off the platform. \\
      Seller mode
        & Sells only its own product and sets price $p_m$.
        & \textemdash\textemdash
        & Provide the outside competitive benchmark. \\
      Dual mode
        & Hosts the marketplace, sets commission $\tau$, sells its own product, and sets price $p_m$.
        & Joins the marketplace, chooses innovation $\Delta$, inside price $p_i$, and outside price $p_o$.
        & Provide the competitive benchmark on and off the platform. \\
      \bottomrule
    \end{tabularx}
  \end{table}

  \vspace{-0.05cm}
  \begin{block}{Equilibrium objects}
    \scriptsize
    For each mode, the analysis pins down prices, seller innovation,
    transaction volume, and platform profit.
  \end{block}
\end{frame}

\begin{frame}
  \frametitle{Marketplace Mode: Setup}
  \footnotesize
  In marketplace mode, $M$ operates as a pure intermediary: it hosts sellers,
  charges a commission $\tau$, and does not sell its own product.

  After joining the marketplace and choosing innovation $\Delta$, seller $S$
  chooses an inside price $p_i$ and an outside price $p_o$.

  \begin{center}
    \hypertarget{mp_price_bounds_src}{%
      \hyperlink{mp_price_bounds_deriv}{%
        $\displaystyle p_i \leq \tau+\Delta, \qquad
        p_o \leq \tau+\Delta-b$}}
  \end{center}

  \begin{block}{Two pricing strategies for $S$}
    \begin{itemize}
      \item \textbf{Sell through $M$:} set $p_o>p_i-b$ and choose
      \[
        \max_{p_i\leq \tau+\Delta}(p_i-\tau)
        G(\nu+b+\Delta-p_i).
      \]
      \item \textbf{Sell directly:} set $p_i>p_o+b$ and choose
      \[
        \max_{p_o\leq \tau+\Delta-b}p_oG(\nu+\Delta-p_o).
      \]
    \end{itemize}
  \end{block}

  \begin{alertblock}{Intuition}
    \textbf{$S$ sells through the platform iff $\tau\leq b$.} If $\tau>b$,
    $S$ uses the platform for discovery but induces consumers to buy directly.
  \end{alertblock}
\end{frame}

\begin{frame}
  \frametitle{Marketplace Equilibrium}
  \footnotesize
  Proposition 1 follows from two binding constraints. Since
  \textbf{$S$ sells through the platform iff $\tau\leq b$}, $M$ faces the
  \textit{showrooming constraint} when choosing its commission.

  \begin{columns}[T,onlytextwidth]
    \begin{column}{0.48\textwidth}
      \begin{block}{\textit{Showrooming constraint}}
        The platform fee binds at
        \begin{center}
          \hypertarget{mp_showrooming_src}{%
            \hyperlink{mp_showrooming_deriv}{%
              $\displaystyle \tau^{mkt}=b$}}
        \end{center}
        $M$ charges the highest fee that still keeps trade on the platform.
      \end{block}
    \end{column}

    \begin{column}{0.48\textwidth}
      \begin{block}{\textit{Competitive constraint}}
        The inside price binds at
        \begin{center}
          \hypertarget{mp_competitive_src}{%
            \hyperlink{mp_competitive_deriv}{%
              $\displaystyle p_i^*=\tau^{mkt}+\Delta^{mkt}
              =b+\Delta^{mkt}$}}
        \end{center}
        Hence demand is
        \[
          G(\nu+b+\Delta^{mkt}-p_i^*)=G(\nu).
        \]
      \end{block}
    \end{column}
  \end{columns}

  \vspace{0.1cm}
  \begin{equation*}
    K'(\Delta^{mkt})=G(\nu), \qquad
    \Pi^{mkt}=bG(\nu), \qquad
    \pi^{mkt}=\Delta^{mkt}G(\nu)-K(\Delta^{mkt}).
  \end{equation*}

  \begin{block}{Intuition}
    $S$ sells exclusively through the marketplace; innovation is chosen until
    marginal innovation cost equals the induced transaction volume.
  \end{block}
\end{frame}

\begin{frame}
  \frametitle{Seller Mode: Setup}
  \footnotesize
  In seller mode, $M$ sells only its own product. Seller $S$ is unavailable on
  the platform, so consumers never discover $S$ through $M$.

  \begin{columns}[T,onlytextwidth]
    \begin{column}{0.48\textwidth}
      \begin{block}{Seller $S$}
        \begin{itemize}
          \item Its outside price is payoff-irrelevant.
          \item It chooses the lowest innovation level:
          \[
            \Delta^{sell}=\Delta^l.
          \]
          \item It makes no sales in equilibrium.
        \end{itemize}
      \end{block}
    \end{column}

    \begin{column}{0.48\textwidth}
      \begin{block}{Platform $M$}
        $M$ chooses the price of its own product:
        \begin{equation}
          \max_{p_m\leq b+\sigma}
          p_m G(\nu+b+\sigma-p_m)
          \label{eq:seller-profit}
        \end{equation}
      \end{block}
    \end{column}
  \end{columns}

  \begin{alertblock}{Price cap}
    The cap $p_m\leq b+\sigma$ comes from competition with fringe sellers in
    the direct channel: $\nu+\sigma+b-p_m\geq\nu$.
  \end{alertblock}
\end{frame}

\begin{frame}
  \frametitle{Seller Mode Equilibrium}
  \footnotesize
  The \textit{competitive constraint} binds in the seller-mode pricing problem,
  so $M$ charges the highest price that still beats fringe sellers.

  \begin{columns}[T,onlytextwidth]
    \begin{column}{0.48\textwidth}
      \begin{block}{Binding price}
        \begin{center}
          \hypertarget{seller_competitive_src}{%
            \hyperlink{seller_competitive_deriv}{%
              $\displaystyle p_m^*=b+\sigma$}}
        \end{center}
        The price equals convenience value plus $M$'s product advantage.
      \end{block}
    \end{column}

    \begin{column}{0.48\textwidth}
      \begin{block}{Demand and profits}
        Substituting $p_m^*$ gives demand
        \[
          G(\nu+b+\sigma-p_m^*)=G(\nu).
        \]
        Therefore
        \[
          \Pi^{sell}=(b+\sigma)G(\nu),
          \qquad
          \pi^{sell}=0.
        \]
      \end{block}
    \end{column}
  \end{columns}

  \begin{block}{Proposition 2}
    $M$ sets $p_m^*=b+\sigma$ and sells to all consumers, whereas $S$ sells to
    no one.
  \end{block}

  \begin{alertblock}{Intuition}
    Seller mode eliminates on-platform competition from $S$: the platform
    captures convenience value $b$ plus its own product advantage $\sigma$.
  \end{alertblock}
\end{frame}

\begin{frame}
  \frametitle{Dual Mode: Setup}
  \footnotesize
  In dual mode, $M$ is both an intermediary and a seller: it hosts $S$, charges
  commission $\tau$, sells its own product, and sets its own price $p_m$.

  \begin{columns}[T,onlytextwidth]
    \begin{column}{0.48\textwidth}
      \begin{block}{Platform $M$}
        \begin{itemize}
          \item sets commission $\tau$;
          \item sells its own product at price $p_m$;
          \item earns either commission income or own-product revenue.
        \end{itemize}
      \end{block}
    \end{column}

    \begin{column}{0.48\textwidth}
      \begin{block}{Seller $S$}
        \begin{itemize}
          \item chooses innovation $\Delta$;
          \item chooses inside price $p_i$ and outside price $p_o$;
          \item competes with both fringe sellers and $M$'s own product.
        \end{itemize}
      \end{block}
    \end{column}
  \end{columns}

  \begin{alertblock}{New force}
    $M$'s own price becomes an extra cap on $S$'s inside price. This is the
    source of the price-squeeze mechanism.
  \end{alertblock}
\end{frame}

\begin{frame}
  \frametitle{Dual Mode: Pricing Subgame}
  \footnotesize
  Given $\tau$ and $\Delta$, the stage-3 pricing outcome depends on whether
  $M$'s own product or $S$'s product has the higher quality.

  \begin{columns}[T,onlytextwidth]
    \begin{column}{0.48\textwidth}
      \begin{block}{Semi-seller equilibrium}
        If $\sigma\geq\Delta$, $M$'s product is weakly better:
        \begin{center}
          \hypertarget{dual_semiseller_src}{%
            \hyperlink{dual_semiseller_deriv}{%
              $\displaystyle p_i^*=\tau,\quad
              p_m^*=\tau+\sigma-\Delta$}}
        \end{center}
        Consumers buy from $M$; $S$ earns zero operating margin.
      \end{block}
    \end{column}

    \begin{column}{0.48\textwidth}
      \begin{block}{Price-squeeze equilibrium}
        If $\Delta>\sigma$, $S$'s product is better:
        \begin{center}
          \hypertarget{dual_squeeze_src}{%
            \hyperlink{dual_squeeze_deriv}{%
              $\displaystyle p_i^*=p_m^*+\Delta-\sigma$}}
        \end{center}
        Consumers buy from $S$, but $M$'s own price caps $p_i^*$.
      \end{block}
    \end{column}
  \end{columns}

  \begin{alertblock}{Intuition}
    Dual mode matters when $S$ is better: $M$ may not sell, yet its own product
    disciplines $S$'s price.
  \end{alertblock}
\end{frame}

\begin{frame}
  \frametitle{Price Squeeze}
  \footnotesize
  When $\Delta>\sigma$, consumers buy from $S$ through the marketplace, but
  $M$'s own price caps $S$'s inside price.

  \begin{columns}[T,onlytextwidth]
    \begin{column}{0.48\textwidth}
      \begin{block}{Feasible own-product price}
        The price-squeeze interval is
        \[
          \hypertarget{dual_pm_interval_src}{%
            \hyperlink{dual_pm_interval_deriv}{%
              p_m^*\in
              \left[
                \max\{\tau-\Delta+\sigma,0\},
                \tau+\min\{\sigma,0\}
              \right]}}
        \]
        and the selected equilibrium uses the lowest feasible $p_m^*$.
      \end{block}
    \end{column}

    \begin{column}{0.48\textwidth}
      \begin{block}{Platform profit}
        Given $p_i^*=p_m^*+\Delta-\sigma$, demand is
        $G(\nu+\sigma+b-p_m^*)$. Hence
        \begin{center}
          \hypertarget{dual_profit_src}{%
            \hyperlink{dual_profit_deriv}{%
              $\displaystyle \Pi=\tau G(\nu+\sigma+b-p_m^*)$}}
        \end{center}
      \end{block}
    \end{column}
  \end{columns}

  \begin{alertblock}{Intuition}
    A lower $p_m^*$ lowers $S$'s feasible price, expands demand, and raises
    commission income. This is why the selection rule picks
    $p_m^*=\max\{\tau-\Delta+\sigma,0\}$.
  \end{alertblock}
\end{frame}

\begin{frame}
  \frametitle{Stage-2 Innovation Decision}
  \footnotesize
  After selecting the lowest feasible $p_m^*$, $S$'s payoff can be written as
  \begin{center}
    \hypertarget{dual_s_profit_src}{%
      \hyperlink{dual_s_profit_deriv}{%
        $\displaystyle
        \pi(\Delta)=
        \max\{(\Delta-\sigma-\tau)G(\nu+\sigma+b),0\}-K(\Delta)$}}
  \end{center}

  \begin{columns}[T,onlytextwidth]
    \begin{column}{0.48\textwidth}
      \begin{block}{Threshold commission}
        Denote $\bar{\tau}\in(\Delta^l-\sigma,\bar{\Delta}-\sigma)$ by
        \[
          \hypertarget{dual_taubar_src}{%
            \hyperlink{dual_taubar_deriv}{%
              (\bar{\Delta}-\sigma-\bar{\tau})G(\nu+b+\sigma)
              -K(\bar{\Delta})=0.}}
        \]
      \end{block}
    \end{column}

    \begin{column}{0.48\textwidth}
      \begin{block}{Lemma 1}
        \begin{itemize}
          \item If $\tau\leq\bar{\tau}$, then
            $\Delta=\bar{\Delta}$.
          \item If $\tau>\bar{\tau}$, then
            $\Delta=\Delta^l$.
        \end{itemize}
      \end{block}
    \end{column}
  \end{columns}

  \begin{alertblock}{Intuition}
    $\bar{\tau}$ is the highest commission compatible with high innovation;
    above it, price squeeze removes $S$'s incentive to pay $K(\bar{\Delta})$.
  \end{alertblock}
\end{frame}

\begin{frame}
  \frametitle{Stage-1 Commission Decision}
  \footnotesize
  In stage 1, $M$ chooses $\tau$ while anticipating the innovation response.
  Two constraints shape the choice.

  \begin{columns}[T,onlytextwidth]
    \begin{column}{0.48\textwidth}
      \begin{block}{Low fee, high innovation}
        Set $\tau\leq\bar{\tau}$, induce $\Delta=\bar{\Delta}$, and earn
        commission on high demand:
        \[
          \tau G(\nu+\sigma+b).
        \]
        At the best such policy, $\tau=\min\{b,\bar{\tau}\}$.
      \end{block}
    \end{column}

    \begin{column}{0.48\textwidth}
      \begin{block}{High fee, low innovation}
        Set the showrooming upper bound $\tau=b$, accept
        $\Delta=\Delta^l$, and earn
        \[
          \bigl(b+\max\{\sigma-\Delta^l,0\}\bigr)G(\nu+\Delta^l).
        \]
      \end{block}
    \end{column}
  \end{columns}

  \begin{alertblock}{Where the constraints come from}
    The \textit{showrooming constraint} $\tau\leq b$ keeps transactions on
    the platform; the \textit{innovation constraint} $\tau\leq\bar{\tau}$
    keeps $S$ willing to choose high innovation.
  \end{alertblock}
\end{frame}

\begin{frame}
  \frametitle{Dual Mode Equilibrium}
  \footnotesize
  Proposition 3 gives $M$'s optimal commission after anticipating $S$'s
  innovation response.

  \begin{itemize}
    \item \textbf{Case 1: $b\leq\bar{\tau}$.} The showrooming upper bound
      already satisfies the innovation constraint:
      \[
        \tau^{dual}=b,\qquad
        \Delta^{dual}=\bar{\Delta},\qquad
        \Pi^{dual}=bG(\nu+\sigma+b).
      \]

    \item \textbf{Case 2: $b>\bar{\tau}$.} $M$ cannot set the high fee $b$
      and preserve high innovation. It compares
      \[
        \hypertarget{dual_prop3_src}{%
          \hyperlink{dual_prop3_deriv}{%
            \bar{\tau}G(\nu+\sigma+b)
            \quad\text{and}\quad
            \bigl(b+\max\{\sigma-\Delta^l,0\}\bigr)G(\nu+\Delta^l).}}
      \]
      If the first term is larger, $M$ sets
      $(\tau,\Delta)=(\bar{\tau},\bar{\Delta})$; otherwise it sets
      $(\tau,\Delta)=(b,\Delta^l)$.

    \item \textbf{Trade-off in Case 2.} A lower fee preserves innovation and
      expands transactions; a higher fee extracts more per transaction but
      may destroy innovation.
  \end{itemize}
\end{frame}

\begin{frame}
  \frametitle{Partial and Complete Price Squeeze}
  \footnotesize
  Price squeeze can be partial or complete depending on whether $S$ keeps a
  positive operating margin after $M$'s selected price.

  \begin{itemize}
    \item \textbf{Partial price squeeze.} With high innovation,
      \[
        \Delta^{dual}=\bar{\Delta},\qquad
        p_m^*=0,\qquad
        p_i^*=\bar{\Delta}-\sigma.
      \]
      If $p_i^*>\tau^{dual}$, then $S$ still earns positive margin:
      \[
        p_i^*-\tau^{dual}>0.
      \]

    \item \textbf{Complete price squeeze.} With low innovation and
      $\Delta^l>\sigma$, $S$ still sells, but $M$ sets $p_m^*$ so that
      \[
        p_i^*=\tau^{dual}.
      \]
      Then $S$ sells with zero operating margin.

    \item \textbf{Boundary case.} If $\Delta^l\leq\sigma$, $M$'s product is
      weakly better and the pricing subgame becomes the semi-seller
      equilibrium.
  \end{itemize}
\end{frame}

\begin{frame}
  \frametitle{Which Mode Does the Platform Choose?}
  \footnotesize
  Corollary 1 compares platform profits across the three operating modes.

  \begin{columns}[T,onlytextwidth]
    \begin{column}{0.48\textwidth}
      \begin{block}{Pure modes}
        \[
          \Pi^{mkt}=bG(\nu),
          \qquad
          \Pi^{sell}=(b+\sigma)G(\nu).
        \]
        Therefore
        \[
          \Pi^{mkt}\geq\Pi^{sell}
          \quad\Longleftrightarrow\quad
          \sigma\leq0.
        \]
      \end{block}
    \end{column}

    \begin{column}{0.48\textwidth}
      \begin{block}{Full choice}
        \begin{itemize}
          \item Dual mode strictly dominates marketplace mode:
            $\Pi^{dual}>\Pi^{mkt}$.
          \item There is a unique $\underline{\sigma}>0$ such that
            \[
              \Pi^{dual}>\Pi^{sell}
              \quad\Longleftrightarrow\quad
              \sigma<\underline{\sigma}.
            \]
          \item Therefore:
            \[
              M \text{ chooses }
              \begin{cases}
                \text{dual mode}, & \sigma<\underline{\sigma},\\
                \text{seller mode}, & \sigma\geq\underline{\sigma}.
              \end{cases}
            \]
        \end{itemize}
      \end{block}
    \end{column}
  \end{columns}

  \begin{alertblock}{Intuition}
    Thus $M$ chooses dual mode when $\sigma<\underline{\sigma}$ and seller
    mode when $\sigma\geq\underline{\sigma}$. For small $\sigma$, price
    squeeze creates more transactions; for large $\sigma$, seller mode lets
    $M$ exploit its own product advantage without revealing $S$.
  \end{alertblock}
\end{frame}

\begin{frame}
  \frametitle{Baseline Ban on Dual Mode: Mode Switch}
  \footnotesize
  Proposition 4 asks what happens if regulation removes dual mode from $M$'s
  choice set.

  \begin{itemize}
    \item \textbf{Post-ban choice set:}
      \[
        \{\text{marketplace},\text{seller},\text{dual}\}
        \quad\longrightarrow\quad
        \{\text{marketplace},\text{seller}\}.
      \]

    \item \textbf{Post-ban mode choice:}
      \[
        M \text{ chooses seller mode iff } \sigma>0,
        \qquad
        M \text{ chooses marketplace mode iff } \sigma\leq0.
      \]

    \item \textbf{Compared with the unconstrained baseline:}
      \[
        \begin{array}{lll}
          \sigma\geq\underline{\sigma}: & \text{seller}\to\text{seller},
          & \text{no real effect},\\
          0<\sigma<\underline{\sigma}: & \text{dual}\to\text{seller},
          & S \text{ disappears from the platform},\\
          \sigma\leq0: & \text{dual}\to\text{marketplace},
          & S \text{ remains but price squeeze disappears}.
        \end{array}
      \]
  \end{itemize}
\end{frame}

\begin{frame}
  \frametitle{Baseline Ban on Dual Mode: Effects}
  \scriptsize

  \begin{table}
    \centering
    \renewcommand{\arraystretch}{1.22}
    \begin{tabularx}{\textwidth}{@{}p{2.4cm}p{2.1cm}p{1.1cm}p{1.1cm}p{1.1cm}p{1.1cm}X@{}}
      \toprule
      Region & Mode switch & $\Pi$ & $\pi$ & $CS$ & $\Delta$ & Welfare logic \\
      \midrule
      $\sigma\geq\underline{\sigma}$
        & seller $\to$ seller
        & $=$ & $=$ & $=$ & $=$
        & Ban is nonbinding because $M$ already prefers seller mode. \\
      $0<\sigma<\underline{\sigma}$
        & dual $\to$ seller
        & $\downarrow$ & $\downarrow$ or $=$ & $\downarrow$ & $\downarrow$ or $=$
        & Welfare falls: the ban removes $S$ from consumer discovery and weakens transactions. \\
      $\sigma\leq0$
        & dual $\to$ marketplace
        & $\downarrow$ & $\uparrow$ & $\downarrow$ & $\downarrow$ or $\uparrow$
        & Welfare is ambiguous; it rises only if innovation gains are large enough. \\
      \bottomrule
    \end{tabularx}
  \end{table}

  \begin{alertblock}{Key implication}
    A dual-mode ban always weakly lowers consumer surplus. It improves welfare
    only in the $\sigma\leq0$ region and only when marketplace innovation
    gains dominate the loss of price-squeeze-driven transaction surplus.
  \end{alertblock}
\end{frame}

\begin{frame}
  \frametitle{When Can a Dual-Mode Ban Raise Welfare?}
  \footnotesize
  In the $\sigma\leq0$ region, the ban changes the market from dual mode to
  marketplace mode. Welfare can rise only if two conditions hold.

  \begin{itemize}
    \item \textbf{Dual mode must have low innovation:}
      \[
        \Delta^{dual}=\Delta^l.
      \]
      Otherwise the ban replaces high dual-mode innovation
      $\bar{\Delta}$ with lower marketplace innovation $\Delta^{mkt}$.

    \item \textbf{Marketplace innovation gains must be large enough:}
      \begin{center}
        \hypertarget{ban_welfare_condition_src}{%
          \hyperlink{ban_welfare_condition_deriv}{%
            $\displaystyle
            (b+\Delta^{mkt})G(\nu)-K(\Delta^{mkt})
            >
            bG(\nu+\Delta^l)
            +
            \int_0^{\nu+\Delta^l}
            \left[
              \nu+\Delta^l-\max\{\nu,\nu_o\}
            \right]dG(\nu_o)$}}
      \end{center}
  \end{itemize}

  \begin{alertblock}{Interpretation}
    The left side is marketplace value from platform convenience plus higher
    innovation. The right side is the dual-mode low-innovation welfare
    benchmark.
  \end{alertblock}
\end{frame}

\begin{frame}
  \frametitle{Derivation: Marketplace Price Bounds}
  \hypertarget{mp_price_bounds_deriv}{}
  \footnotesize

  \begin{block}{On-platform price bound}
    Fringe sellers on $M$ provide utility $\nu+b-\tau$. If consumers buy
    $S$ on $M$, utility is $\nu+\Delta+b-p_i$. Hence $S$ must satisfy
    \[
      \nu+\Delta+b-p_i \geq \nu+b-\tau
      \quad\Longrightarrow\quad
      p_i\leq \tau+\Delta.
    \]
  \end{block}

  \begin{block}{Direct-channel price bound}
    If consumers buy directly from $S$, utility is $\nu+\Delta-p_o$.
    Comparing this with fringe sellers on $M$ gives
    \[
      \nu+\Delta-p_o \geq \nu+b-\tau
      \quad\Longrightarrow\quad
      p_o\leq \tau+\Delta-b.
    \]
  \end{block}

  \hyperlink{mp_price_bounds_src}{\beamerbutton{Back to price bounds}}
\end{frame}

\begin{frame}
  \frametitle{Derivation: Showrooming Constraint}
  \hypertarget{mp_showrooming_deriv}{}
  \footnotesize

  \begin{block}{Compare $S$'s two channels}
    Under the binding price bounds, both channels induce the same transaction
    utility, $\nu+b-\tau$, but they give $S$ different margins:
    \[
      \text{through }M:\;(\tau+\Delta)-\tau=\Delta,
      \qquad
      \text{directly}:\;\tau+\Delta-b=\Delta+\tau-b.
    \]
  \end{block}

  \begin{alertblock}{Channel choice}
    $S$ sells through $M$ iff
    \[
      \Delta \geq \Delta+\tau-b
      \quad\Longleftrightarrow\quad
      \tau\leq b.
    \]
    Therefore any $\tau>b$ makes $S$ use $M$ for discovery but shift
    transactions to the direct channel.
  \end{alertblock}

  \hyperlink{mp_showrooming_src}{\beamerbutton{Back to showrooming formula}}
\end{frame}

\begin{frame}
  \frametitle{Derivation: Marketplace Competitive Constraint}
  \hypertarget{mp_competitive_deriv}{}
  \footnotesize

  For inside sales, $S$ solves
  \[
    \max_{p_i\leq\tau+\Delta}(p_i-\tau)G(\nu+b+\Delta-p_i).
  \]

  \begin{block}{Why the upper bound binds}
    The derivative at the competitive upper bound $p_i=\tau+\Delta$ is
    \[
      G(\nu+b-\tau)-\Delta g(\nu+b-\tau).
    \]
    Since $\tau\leq b$ and $G/g$ is increasing under log-concavity of $G$,
    Assumption (2) implies
    \[
      \Delta\leq\bar{\Delta}<\frac{G(\nu)}{g(\nu)}
      \leq
      \frac{G(\nu+b-\tau)}{g(\nu+b-\tau)}.
    \]
    The derivative is positive at the bound, so $S$ would like to raise
    price further; the \textit{competitive constraint} stops it.
  \end{block}

  \hyperlink{mp_competitive_src}{\beamerbutton{Back to competitive formula}}
\end{frame}

\begin{frame}
  \frametitle{Derivation: Seller Competitive Constraint}
  \hypertarget{seller_competitive_deriv}{}
  \footnotesize

  In seller mode, $M$ solves
  \[
    \max_{p_m\leq b+\sigma}p_mG(\nu+b+\sigma-p_m).
  \]

  \begin{block}{Why the upper bound binds}
    At the competitive upper bound $p_m=b+\sigma$, the derivative is
    \[
      G(\nu)-(b+\sigma)g(\nu).
    \]
    Assumption (2) gives
    \[
      b+\sigma<\frac{G(\nu)}{g(\nu)},
    \]
    so the derivative is positive at the bound. Therefore $M$ would like
    to raise price further, but the direct-channel fringe option imposes
    the \textit{competitive constraint}.
  \end{block}

  \hyperlink{seller_competitive_src}{\beamerbutton{Back to seller equilibrium}}
\end{frame}

\begin{frame}
  \frametitle{Derivation: Semi-Seller Pricing}
  \hypertarget{dual_semiseller_deriv}{}
  \footnotesize

  \begin{block}{Why $S$ prices at effective marginal cost}
    When $\sigma\geq\Delta$, $M$'s own product is weakly better. Seller $S$
    pays commission $\tau$ on each platform sale, so its effective marginal
    cost is $\tau$:
    \[
      p_i^*=\tau.
    \]
  \end{block}

  \begin{block}{Why $M$ sells at $p_m^*=\tau+\sigma-\Delta$}
    Consumers are indifferent between $M$ and $S$ when
    \[
      \nu+\sigma+b-p_m^*
      =
      \nu+\Delta+b-p_i^*.
    \]
    Substituting $p_i^*=\tau$ gives
    \[
      p_m^*=\tau+\sigma-\Delta.
    \]
  \end{block}

  \hyperlink{dual_semiseller_src}{\beamerbutton{Back to pricing cases}}
\end{frame}

\begin{frame}
  \frametitle{Derivation: Price-Squeeze Price Cap}
  \hypertarget{dual_squeeze_deriv}{}
  \footnotesize

  \begin{block}{Price cap imposed by $M$}
    Consumers buy $S$ rather than $M$ only if
    \[
      \nu+\Delta+b-p_i
      \geq
      \nu+\sigma+b-p_m,
      \quad\Longrightarrow\quad
      p_i\leq p_m+\Delta-\sigma.
    \]
    The competitive constraint binds, so $p_i^*=p_m^*+\Delta-\sigma$.
  \end{block}

  \begin{alertblock}{Why this is price squeeze}
    $M$ may not sell in equilibrium, but a lower $p_m^*$ tightens the upper
    bound on $S$'s price and forces $S$ to charge less.
  \end{alertblock}

  \hyperlink{dual_squeeze_src}{\beamerbutton{Back to price squeeze relation}}
\end{frame}

\begin{frame}
  \frametitle{Derivation: Feasible Interval for $p_m^*$}
  \hypertarget{dual_pm_interval_deriv}{}
  \scriptsize

  \begin{columns}[T,onlytextwidth]
    \begin{column}{0.48\textwidth}
      \begin{block}{Upper bounds}
        \begin{itemize}
          \item $p_m^*\leq\tau$: otherwise $M$ could undercut $S$ and earn
            an own-product margin above the commission $\tau$.
          \item $p_m^*\leq\tau+\sigma$: otherwise consumers prefer platform
            fringe sellers, since
            \[
              \nu+\sigma+b-p_m^*<\nu+b-\tau.
            \]
        \end{itemize}
        Hence $p_m^*\leq\tau+\min\{\sigma,0\}$.
      \end{block}
    \end{column}

    \begin{column}{0.48\textwidth}
      \begin{block}{Lower bounds}
        \begin{itemize}
          \item $S$ must not sell below effective marginal cost:
            $p_i^*\geq\tau$.
          \item Since $p_i^*=p_m^*+\Delta-\sigma$,
            $p_m^*\geq\tau-\Delta+\sigma$.
          \item $M$'s own price is nonnegative: $p_m^*\geq0$.
        \end{itemize}
        Hence $p_m^*\geq\max\{\tau-\Delta+\sigma,0\}$.
      \end{block}
    \end{column}
  \end{columns}

  \vspace{-0.05cm}
  \[
    p_m^*\in
    \left[
      \max\{\tau-\Delta+\sigma,0\},
      \tau+\min\{\sigma,0\}
    \right].
  \]

  \hyperlink{dual_pm_interval_src}{\beamerbutton{Back to feasible interval}}
\end{frame}

\begin{frame}
  \frametitle{Derivation: Selected Price in Price Squeeze}
  \footnotesize

  In a price-squeeze equilibrium, consumers buy from $S$ and $M$ earns
  commission revenue:
  \[
    \Pi=\tau G(\nu+\sigma+b-p_m^*).
  \]

  \begin{block}{Selection rule}
    Because $G(\cdot)$ is increasing, a lower $p_m^*$ raises transaction
    volume and therefore raises $M$'s commission profit. The selected
    equilibrium sets
    \[
      p_m^*=\max\{\tau-\Delta+\sigma,0\}.
    \]
  \end{block}

  \hyperlink{dual_pm_interval_src}{\beamerbutton{Back to selected price}}
\end{frame}

\begin{frame}
  \frametitle{Derivation: Platform Profit in Price Squeeze}
  \hypertarget{dual_profit_deriv}{}
  \footnotesize

  In the price-squeeze equilibrium, consumers buy $S$ through $M$, so the
  platform earns commission $\tau$ per transaction.

  \begin{block}{Demand under price squeeze}
    Using $p_i^*=p_m^*+\Delta-\sigma$,
    \[
      \nu+\Delta+b-p_i^*
      =
      \nu+\Delta+b-(p_m^*+\Delta-\sigma)
      =
      \nu+\sigma+b-p_m^*.
    \]
    Demand is therefore $G(\nu+\sigma+b-p_m^*)$.
  \end{block}

  \begin{block}{Commission profit}
    \[
      \Pi
      =
      \tau G(\nu+\sigma+b-p_m^*).
    \]
    Since $G(\cdot)$ is increasing, the selected lowest $p_m^*$ maximizes
    $M$'s commission profit among price-squeeze equilibria.
  \end{block}

  \hyperlink{dual_profit_src}{\beamerbutton{Back to platform profit}}
\end{frame}

\begin{frame}
  \frametitle{Derivation: Stage-2 Innovation Decision}
  \hypertarget{dual_s_profit_deriv}{}
  \hypertarget{dual_taubar_deriv}{}
  \footnotesize

  With the selected lowest $p_m^*$, $S$'s payoff combines the zero-margin
  semi-seller case and the positive-margin price-squeeze case:
  \[
    \pi(\Delta)=
    \max\{(\Delta-\sigma-\tau)G(\nu+\sigma+b),0\}-K(\Delta).
  \]

  \begin{block}{Definition of the innovation threshold}
    Lemma 1 denotes
    $\bar{\tau}\in(\Delta^l-\sigma,\bar{\Delta}-\sigma)$ as the commission
    that makes $S$ indifferent between high innovation and the zero-profit
    low-innovation option:
    \[
      (\bar{\Delta}-\sigma-\bar{\tau})G(\nu+b+\sigma)
      -K(\bar{\Delta})=0.
    \]
    Equivalently,
    \[
      \bar{\tau}
      =
      \bar{\Delta}-\sigma
      -
      \frac{K(\bar{\Delta})}{G(\nu+b+\sigma)}.
    \]
  \end{block}

  \hyperlink{dual_s_profit_src}{\beamerbutton{Back to seller payoff}}
  \hyperlink{dual_taubar_src}{\beamerbutton{Back to innovation threshold}}
\end{frame}

\begin{frame}
  \frametitle{Derivation: Proposition 3 Profit Comparison}
  \hypertarget{dual_prop3_deriv}{}
  \scriptsize

  If $b\leq\bar{\tau}$, the showrooming upper bound also satisfies the
  innovation constraint, so $M$ sets $\tau^{dual}=b$ and induces
  $\Delta^{dual}=\bar{\Delta}$.

  \vspace{0.05cm}
  If $b>\bar{\tau}$, $M$ cannot set $\tau=b$ and preserve high innovation.
  It compares two feasible policies.

  \begin{columns}[T,onlytextwidth]
    \begin{column}{0.48\textwidth}
      \begin{block}{Low fee, high innovation}
        Set $\tau=\bar{\tau}$ and induce $\Delta=\bar{\Delta}$:
        \[
          \Pi_H
          =
          \bar{\tau}G(\nu+\sigma+b).
        \]
      \end{block}
    \end{column}

    \begin{column}{0.48\textwidth}
      \begin{block}{High fee, low innovation}
        Set $\tau=b$ and accept $\Delta=\Delta^l$:
        \[
          \Pi_L
          =
          \bigl(b+\max\{\sigma-\Delta^l,0\}\bigr)
          G(\nu+\Delta^l).
        \]
      \end{block}
    \end{column}
  \end{columns}

  \begin{alertblock}{Condition (8)}
    $M$ preserves high innovation when
    \[
      \bar{\tau}G(\nu+\sigma+b)
      \geq
      \bigl(b+\max\{\sigma-\Delta^l,0\}\bigr)G(\nu+\Delta^l).
    \]
    Otherwise it sets $\tau^{dual}=b$ and accepts
    $\Delta^{dual}=\Delta^l$.
  \end{alertblock}

  \hyperlink{dual_prop3_src}{\beamerbutton{Back to Proposition 3}}
\end{frame}

\begin{frame}
  \frametitle{Derivation: Welfare Condition Under Dual-Mode Ban}
  \hypertarget{ban_welfare_condition_deriv}{}
  \scriptsize

  The welfare-improving case requires $\sigma\leq0$ and
  $\Delta^{dual}=\Delta^l$. The ban then replaces low-innovation dual mode
  with marketplace mode.

  \begin{block}{Marketplace welfare component}
    In marketplace mode, the binding price is $p_i^*=b+\Delta^{mkt}$ and
    $S$ chooses $K'(\Delta^{mkt})=G(\nu)$. The relevant surplus generated by
    platform convenience and innovation is
    \[
      (b+\Delta^{mkt})G(\nu)-K(\Delta^{mkt}).
    \]
  \end{block}

  \begin{block}{Low-innovation dual-mode benchmark}
    With $\Delta^{dual}=\Delta^l$, price squeeze generates transaction value
    $bG(\nu+\Delta^l)$ plus consumer gains from the default innovation:
    \[
      \int_0^{\nu+\Delta^l}
      \left[
        \nu+\Delta^l-\max\{\nu,\nu_o\}
      \right]dG(\nu_o).
    \]
  \end{block}

  \begin{alertblock}{Condition (9)}
    The ban raises welfare only if the marketplace component exceeds this
    low-innovation dual-mode benchmark.
  \end{alertblock}

  \hyperlink{ban_welfare_condition_src}{\beamerbutton{Back to welfare condition}}
\end{frame}
